
\chapter{Introduction}

\section{Project Information}

This project is developed as part of the Software Engineering course with the objective of applying software engineering principles, methodologies, and tools to the development of a real-world software system.

The project focuses on designing and implementing a centralized platform for managing horse racing tournaments, supporting various stakeholders throughout the tournament lifecycle, from registration and scheduling to result management and ranking updates.

\begin{table}[H]
	\centering
	\caption{Project Information}
	\begin{tabular}{|p{5cm}|p{8cm}|}
		\hline
		\textbf{Item} & \textbf{Description} \\
		\hline
		Project Name & Horse Racing Tournament Management System \\
		\hline
		Project Acronym & HRTMS \\
		\hline
		Project Group & Group 3 \\
		\hline
		Software Type & Web Application \\
		\hline
		Development Methodology & Agile \\
		\hline
		Frontend Technology & React.js / Next.js \\
		\hline
		Backend Technology & FastAPI \\
		\hline
		Database Management System & Microsoft SQL Server \\
		\hline
		Version & MVP Version 1.0 \\
		\hline
		Academic Course & Software Engineering \\
		\hline
		Academic Year & 2025--2026 \\
		\hline
	\end{tabular}
\end{table}

\section{Project Team}

The Horse Racing Tournament Management System is developed by Group 3. The project team consists of seven members who collaboratively participate in project management, requirement analysis, system design, implementation, testing, and documentation activities.

\begin{table}[H]
	\centering
	\caption{Team Members}
	\begin{tabular}{|c|p{6cm}|p{4cm}|p{3cm}|}
		\hline
		\textbf{No.} & \textbf{Full Name} & \textbf{Student ID} & \textbf{Role} \\
		\hline
		1 & ................................ & ........................ & Team Leader \\
		\hline
		2 & ................................ & ........................ & Member \\
		\hline
		3 & ................................ & ........................ & Member \\
		\hline
		4 & ................................ & ........................ & Member \\
		\hline
		5 & ................................ & ........................ & Member \\
		\hline
		6 & ................................ & ........................ & Member \\
		\hline
		7 & ................................ & ........................ & Member \\
		\hline
	\end{tabular}
\end{table}

Each member is assigned specific responsibilities according to their expertise and contributes to the successful completion of the project deliverables.

\section{Project Background}

Horse racing is one of the oldest sports and entertainment activities in the world, involving multiple stakeholders such as horse owners, jockeys, referees, tournament organizers, and spectators. Organizing a horse racing tournament requires managing a large amount of information, including horse profiles, jockey assignments, race schedules, competition results, and ranking standings.

In many cases, these activities are still performed manually using spreadsheets, paper documents, or disconnected software tools. Such approaches often lead to data inconsistency, scheduling conflicts, delayed result updates, and difficulties in tracking tournament information. As the scale of tournaments grows, managing these processes manually becomes increasingly inefficient and error-prone.

The rapid development of information technology provides an opportunity to digitize and automate these management activities. A centralized web-based system can significantly improve operational efficiency, data accuracy, and transparency for all participants involved in the tournament.

Therefore, the development of a Horse Racing Tournament Management System is proposed to support the management and operation of horse racing tournaments through an integrated software platform.

\section{Problem Statement}

The current management process of horse racing tournaments faces several challenges.

First, information related to horses, jockeys, race schedules, and tournament registrations is often stored separately across different documents or systems. This makes it difficult for organizers to maintain data consistency and perform effective monitoring.

Second, the registration and approval process for tournament participation is usually handled manually. Such a process can result in incomplete information, approval mistakes, or delays in communication between stakeholders.

Third, race scheduling is a complex task that requires coordination among horses, jockeys, referees, and race events. Without proper system support, scheduling conflicts may occur, reducing the overall quality of tournament management.

Finally, competition results and ranking standings are often updated manually, making it difficult for participants and spectators to obtain timely and accurate information.

These issues highlight the need for an integrated management system capable of supporting the complete tournament lifecycle, from registration to result publication.

\section{Software Product Vision}

The vision of the Horse Racing Tournament Management System is to provide a centralized, reliable, and user-friendly platform for managing horse racing tournaments efficiently and transparently.

The system aims to connect all stakeholders, including administrators, horse owners, jockeys, referees, and spectators, within a unified environment where tournament information can be managed and accessed in real time.

By digitalizing the tournament management process, the system seeks to reduce manual workload, minimize operational errors, improve data consistency, and enhance communication among participants. Furthermore, the platform establishes a scalable foundation that can support future enhancements such as notification services, analytics dashboards, prediction systems, and advanced tournament management capabilities.

The long-term vision is to transform traditional horse racing tournament administration into a modern digital ecosystem that improves operational efficiency and user experience for all stakeholders.

\section{Project Scope}

The project focuses on developing a Minimum Viable Product (MVP) that demonstrates the complete operational flow of a horse racing tournament.

The system includes the following major functional areas:

\begin{itemize}
	\item Authentication and Authorization
	\item User Management
	\item Horse Management
	\item Jockey Management
	\item Tournament Management
	\item Registration Management
	\item Race Scheduling
	\item Result Management
	\item Ranking Management
	\item Spectator View
\end{itemize}

Several advanced features are intentionally excluded from the current project scope, including online payment processing, real-money betting, live video streaming, mobile native applications, IoT-based horse tracking, and AI-powered race prediction. These features may be considered for future development phases.

\section{Development Methodology}

The project is developed following Agile Software Development principles. Agile enables the team to divide the project into smaller deliverable components, continuously review progress, and adapt to changing requirements throughout the development process.

Project management activities are conducted using Jira with a Kanban workflow. Development tasks are organized into Epics and Tasks corresponding to major system modules. Source code management is performed using Git and GitHub, following a feature-branch strategy to facilitate collaboration among team members.

The development process consists of requirement analysis, system design, implementation, testing, and final deployment preparation. UML diagrams, database design artifacts, and project documentation are produced throughout the project lifecycle to ensure consistency and maintainability.

\section{Report Structure}

This report is organized into several chapters.

\begin{itemize}
	\item Chapter 1 introduces the project background, objectives, scope, and development methodology.
	\item Chapter 2 presents the software requirements analysis, including stakeholders, functional requirements, and use case analysis.
	\item Chapter 3 describes the system design, architecture, database design, and UML diagrams.
	\item Chapter 4 discusses system implementation details, including frontend, backend, and database development.
	\item Chapter 5 presents testing activities and evaluation results.
	\item Chapter 6 concludes the project and discusses future improvements.
\end{itemize}
